\section{Preserving a cheap observer after marked-work exhaustion}
\label{sec:shadow-fallback}

Block products reduce the cost of many determinant queries, but a large
irreducible cofactor can still exhaust its optional allowance.  The earlier
scheduler then stopped \emph{all} observations, including cheap component-Euler
checks.  We now distinguish those resource limits.  This is analogous to introsort
at the scheduling level: an expensive optional method
has a bounded attempt, then gives way to a cheaper exact method, then to the
complete scanner when its own allowance is exhausted.

The marked backend has the one-way progression
\[
 \text{marked residues and Euler}
 \ \longrightarrow\ \text{Euler only}
 \ \longrightarrow\ \text{unobserved capped scan}.
\]
The first transition occurs on marked-work exhaustion; shared query
exhaustion can take either observer state directly to the unobserved scan.
The crossing order, simplified relative complex, completed pivots, exact
cached values, and global deadline survive every transition.  There is no
restart and no new query allowance.

\subsection{Two limits with separate meanings}

Let $W$ be \code{shadow\_max\_work} and $K$ be
\code{euler\_max\_states}.  While marked observations are active, $W$ pays
for their geometry, matrix allocation, arithmetic updates, and accompanying
Euler queries, as before.  The exception \code{ShadowWorkBudget}, a subtype
of \code{EulerBudget}, identifies exhaustion of this particular allowance.
The driver then stops requesting marked observations permanently.

If query capacity remains, later Euler queries use the direct closure
connectivity evaluator with the \emph{same} exact cache and state counter.
Its new geometric work is separately measured by
\code{euler\_fallback\_traversed\_darts}; it is not charged to, or hidden by
resetting, the already spent marked-work allowance.  The corresponding
\code{euler\_fallback\_evaluations} and \code{euler\_fallback\_states}
fields distinguish all completed fallback requests from newly computed
values.  Cached values do not consume another query, but still check the
global deadline.  Result fields \code{shadow\_exhausted} and
\code{euler\_exhausted}, together with \code{observer\_mode}, distinguish
which methods remain available.

This is an explicit resource-semantics change: setting $W=0$ disables marked
work but permits the query-capped Euler reserve.  Setting $K=0$ disables all
new completion queries.  Either limit may be explicitly unbounded through
the existing API option \code{None}; an unbounded $K$ does not give a finite
fallback-query guarantee.  No additional public switch is required.

If $k$ distinct Euler values have already been computed, the fallback can
compute at most $K-k$ more.  Each direct completion walk uses $O(n)$
combinatorial work, so additional \emph{new closure geometry} is
$O(n(K-k))$.  This bound does not include every observation cost: aggregating
whole relative components still visits the current generators, and cache
keys encode boundary matchings.  Integer bit costs also depend on the
suffix size and coefficient growth.  In particular, this is not a
polynomial bound on the complete Khovanov scan, whose number of live
generators remains uncontrolled in general.

\subsection{Why partially completed observations are safe to abandon}

For a whole relative component, completion gives a genuine chain complex
and its signed Euler characteristic $\chi$.  Therefore
$|\chi|\le\dim H$.  Summing these inequalities with the already certified
nonnegative component multiplicities gives the existing Euler rank lower
bound.  Capping the final nonnegative bound preserves its decision meaning;
the individual signed Euler values themselves remain exact integers.
Thus an Euler certificate does not depend on finishing a marked determinant.

The implementation first completes the component-Euler bound for a stage.
If marked work fails afterward, the incomplete marked bound is discarded.
That stage's already completed Euler bound remains valid; later stages may
request further Euler values.  If the marked allowance instead runs out
while preparing the Euler part, its adapter switches before completing the
outstanding query.  Both paths use the same budget-checked evaluator and
cache.  A generic query-cap exception stops all further observations.
Global \code{ScanLimit} exceptions are never converted into local declines:
they propagate to \code{UNKNOWN} at the recognition boundary.

As before, an early lower bound can certify knottedness only.  An unknot
verdict requires the completed rank calculation.  These invariants hold
for every allowed transition regardless of how favorable the timing turns
out to be.  The policy provides bounded optional work and exact continuation;
it supplies no general quasi-polynomial recognition theorem.

\subsection{Transition tests and benchmark design}

Five new tests exercise zero marked allowance, exhaustion during a block
determinant, one versus two available queries, preserved cached Euler
values, and a global deadline expiring during the transition.  Two hundred
budgeted scans of twenty small random diagrams and their mirrors agree with
independent reduced homology.  Zero-cap and low-cap paths are checked for
both knotted and unknotted outcomes, and no partial observation returns an
unknot verdict.  Existing fallback expectations now reflect the retained
Euler phase.  All 582 maintained tests pass; logs are
\path{data/shadow-fallback-tests.txt} and
\path{data/shadow-fallback-integrated-tests.txt}.

\path{fast/benchmark_shadow_fallback.py} compares the earlier scheduler from
revision \code{82fdfe15b} with the new one, sharing all other dependencies.
It times seven shuffled paired rounds with an identical baseline control,
fresh input and observer setup, and one excluded warmup.  Every successful
arm must agree on the verdict.  Method and stage differences are retained
as evidence, and any \code{UNKNOWN} suppresses all ratios for that case.
The raw backend and full recognition are separate scopes.

Alongside all sixteen existing examples under default limits, the benchmark
uses reduced-work and single-query controls.  A larger default-limit example
joins the Conway knot to $T(3,151)$, giving a validated 313-crossing PD code.
It is tested with the ordinary heuristic order and with a supplied order
that processes the Conway summand first.  The latter case is explicitly
raw-only; its favorable ordering is not attributed to the default heuristic.
Complete input codes, actual scan orders, budgets, source hashes, and raw
samples are retained in
\path{fast/results/shadow_fallback_20261008.json}.

\input{tables/shadow_fallback_benchmark.tex}

\paragraph{Retained query capacity produces the gain.}
On the eightfold Conway sum, $W=0$ improves the raw scan by $5.53$ times,
and $W=2800$ by $5.15$ times.  Both schedulers respect the same marked-work
cap.  In the latter case each spends all 2,800 units and has computed one
Euler value; the new scheduler uses one further query to obtain the
component-Euler certificate at crossing 11.  The baseline completes all 88
crossings.  With $K=1$, neither can compute that further value and both
finish the full scan, with a timing ratio of $1.01$.  This negative control
checks that the improvement does not come from silently resetting the query
budget.  With default marked work, this example never switches and its ratio
is $0.997$.

The retained reserve can also cost time without strengthening a bound.
With $W=0$, $T(3,5)$ uses three fallback values but still requires closed
rank, with ratio $0.982$.  The hard eight-crossing unknot uses one value and
still finishes the entire scan, with ratio $0.984$.  These small regressions
are part of the comparison, not excluded as inconclusive examples.

\paragraph{Ordering determines when the reserve can help.}
For the 313-crossing connected sum, both policies stop marked work at 999,932
units under the default million-unit allowance.  The heuristic order first
processes the torus part.  Fallback gives a certificate at crossing 302
rather than 313, but the measured speedup is only $1.016$, with an A/A ratio
of $1.009$.  This does not demonstrate a substantial timing gain.
With the supplied order, fallback needs only one new query and certifies
knottedness at crossing 11.  Its measured raw-scan speedup is $4.79$ times,
with an A/A ratio of $1.007$.  The respective median times are about
170 and 821 milliseconds.  This is a scheduler gain for the supplied order,
not evidence that the ordinary order selector discovers it.

All 45 case/scope comparisons complete and agree on verdict.  None of the
sixteen default-limit examples switches modes; their raw ratios range from
$0.971$ to $1.031$.  Earlier recognition filters decide every full-recognition
case before this observer runs, giving ratios from $0.982$ to $1.019$ and
no demonstrated full-recognizer gain.  In particular, the larger connected
sum is already detected by the Rasmussen interval filter.  The new policy
therefore improves robustness of an optional backend under its resource
limits while leaving the universal complexity question unresolved.
